\subsection{Motivation for Explicit Autoregressive Mechanisms}
\label{sec:generator-ar-motivation}

Synthetic pretraining motivates a generator that combines diversity with a
retained description of how each series was produced. KernelSynth, introduced
with Chronos, samples sums and products of Gaussian-process kernels and then
draws complete trajectories from the resulting covariance
\citep[Section~4.2]{ansari2024chronos}. Chronos-2 uses it alongside AR, ETS, and
TSI base generators before introducing multivariate dependence
\citep{ansari2025chronos2}. TimePFN's LMC-Synth mixes independent KernelSynth
latent trajectories through a shared linear mixing matrix
\citep{taga2025timepfn}. Toto~2.0 uses the separate \emph{TempoPFN} prior,
which combines several procedural and stochastic families
\citep{moroshan2025tempopfn,khwaja2026toto2}. Section~1.3 describes these
generation procedures and their computational organization.

An explicit autoregressive equation offers a complementary representation:
each dependence is attached to a source series, a lag, and an operation.
For a scalar linear example,
\begin{equation}
 x_t=b+\sum_{\ell=1}^{L}c_{\ell}x_{t-\ell}+\varepsilon_t.
 \label{eq:implemented-ar}
\end{equation}
The nonlinear AST generator replaces the weighted sum with an expression
containing selected lags and interactions. Here $L$ is the largest lag required
by the mechanism; a forecasting model may receive a longer observed context.

The exploratory AR studies motivate this representation by showing strong
one-step approximation for most evaluated KernelSynth and TimePFN latent
series. Section~\ref{sec:ar-coverage-results} presents the measured coverage
and its limitations; the corresponding experimental setup is in
Section~\ref{sec:ar-coverage-protocol}. These studies concern approximation
of observed dynamics, while recovery of the original stochastic mechanism
remains a separate question.
